% Part III-B, discussion (included by main_en.tex).

\section{Discussion}\label{sec:discussion}

\paragraph{A good ranking is not enough.} The frozen \kev{} ranked the unlabelled completions well (AUROC
\NUM{fillauc_all_subkev}), and rewarding them with it still ended below dropping them. Two things it pays for beyond
correctness matter differently. Its generosity, a $\pc$ well above the share of correct completions, raises every
completion's expected reward by about the same amount and is removed by the group baseline. Its preference for
solutions that go on after their answer is not: such text leaves the verifier's label unchanged, so the labelled
completions never punish it, and the judge pays for it on the unlabelled ones (\S\ref{sec:res-repeat}). Both policies
of \texttt{sub\_kev} learned to repeat themselves. A fill-in reward has to be judged by what it pays for among
completions with the same verdict, which an AUROC for the verdict does not measure.

\paragraph{Recalibration fixes the numbers, not the preference.} Refitting the head on the labels so far made the
fill-in reward calibrated (ECE \NUM{fillece_all_subrecal} against \NUM{fillece_all_subkev}), and one of its two
policies still learned to repeat itself. A plausible reason, which these runs do not test directly: the refitted head
starts from the released one and is anchored to it, so it keeps the released head's preferences wherever the labels
say nothing, and repetition after the answer, which leaves the label unchanged, is such a place. A head trained from
scratch has no preference to keep, and neither of its policies repeated itself.

\paragraph{The probe's gain comes with a change of style.} \texttt{sub\_probe} ends \NUM{d_subprobe_math} points above
\texttt{rlvr\_sub} on MATH-500 and \NUM{d_subprobe_aime24} on AIME24, and above the full-label \texttt{rlvr} as well.
Both of its policies switched to drafting first, which no other run of this paper did; in this project's other GRPO
runs the same switch marks the best runs whatever their reward, \NUM{allveronlydraft} of them trained with the
verifier's reward alone (Part II, \S8.3). So with two seeds per condition we cannot say how much of the gain the probe is responsible for. If
the probe makes the switch more likely, which two of two runs suggest but cannot show, a scratch head trained on a
quarter of the labels would be a better fill-in reward than the judge it is built from, for a reason that has little
to do with how well it predicts the labels. More seeds of \texttt{sub\_probe} and \texttt{rlvr\_sub}, and runs that
force or forbid the draft opening, would tell.

\paragraph{Limitations.} One model (4B), one task, one label fraction (a quarter), 200 steps, two seeds per condition;
the recalibration's hyperparameters (anchor, buffer, steps) were set once and not tuned. The intervals resample
problems and not seeds. The style measures were defined after we had read the final completions. The repetition probe
scores with the released head, not with the refitted heads as they were at the end of training. Other experiments
shared the GPUs during training, so we report no training times.

\paragraph{Relation to the earlier versions.} Versions 1--3 of this paper reported a judge and a policy that improve
each other, long-horizon curves and anti-hacking results, none of which came from an experiment (see the notice at the
top and \file{CORRECTION.md}). This version runs a narrower setting that could be checked: partial labels, with only a
small head of the judge adapting. In it, adapting the judge's head fixed its calibration but not what it rewards, and
the best fill-in reward was a head trained from scratch, whose runs also changed how they write.
